\section{Sort a Linked List}
\label{sec:ll-sort-list}

\problemstatement{Given the head of a linked list, sort it in ascending
order and return the new head, ideally in $O(n\log n)$ time and $O(1)$
auxiliary space.}

\begin{patternbox}
\textbf{Merge sort} is the natural sort for linked lists: splitting by the
middle (slow/fast) and merging two sorted lists are both cheap pointer
operations, and unlike arrays no random access is needed. Quicksort's
partition is awkward on lists, so merge sort is the standard choice.
\end{patternbox}

\subsection*{Split, sort halves, merge}

Find the middle with the tortoise--hare and cut the list into two halves.
Recursively sort each half. Then merge the two sorted halves by repeatedly
appending the smaller current node -- exactly the merge step of array merge
sort, but relinking nodes instead of copying values. The recursion bottoms
out at single-node (already sorted) lists.

\begin{ideabox}[Lists Merge and Split Without Extra Arrays]
Splitting a list is just severing it at the middle; merging is relinking
nodes in order. Both are $O(1)$-space pointer work, so merge sort on a list
needs no auxiliary array -- its recursion stack is the only extra space.
\end{ideabox}

\subsection*{Algorithm}

\begin{enumerate}
  \item If the list has $0$ or $1$ nodes, return it.
  \item Split at the middle into two halves.
  \item Recursively sort both halves; merge them into one sorted list.
\end{enumerate}

\subsection*{Dry run}

\dryrun{$3\to1\to2$.}

\begin{itemize}
  \item Split into $3\to1$ and $2$. Sort $3\to1$ into $1\to3$.
  \item Merge $1\to3$ with $2$: $1\to2\to3$.
\end{itemize}

\subsection*{C++ implementation}

\begin{lstlisting}
#include <bits/stdc++.h>
using namespace std;

struct ListNode {
    int val; ListNode* next;
    ListNode(int v) : val(v), next(nullptr) {}
};

ListNode* build(const vector<int>& a) {
    ListNode dummy(0), *t = &dummy;
    for (int x : a) { t->next = new ListNode(x); t = t->next; }
    return dummy.next;
}

ListNode* merge(ListNode* a, ListNode* b) {
    ListNode dummy(0), *t = &dummy;
    while (a && b) {
        if (a->val <= b->val) { t->next = a; a = a->next; }
        else                  { t->next = b; b = b->next; }
        t = t->next;
    }
    t->next = a ? a : b;                          // attach the remainder
    return dummy.next;
}

ListNode* sortList(ListNode* head) {
    if (!head || !head->next) return head;
    ListNode* slow = head; ListNode* fast = head->next;
    while (fast && fast->next) { slow = slow->next; fast = fast->next->next; }
    ListNode* mid = slow->next;
    slow->next = nullptr;                         // cut into two halves
    return merge(sortList(head), sortList(mid));
}

void print(ListNode* h) { for (; h; h = h->next) cout << h->val << " "; cout << "\n"; }

int main() {
    print(sortList(build({4, 2, 1, 3})));         // 1 2 3 4
    return 0;
}
\end{lstlisting}

\begin{complexitybox}
\textbf{Time Complexity.} $O(n\log n)$ -- $\log n$ split levels, $O(n)$
merge per level. \textbf{Space Complexity.} $O(\log n)$ recursion stack.
\end{complexitybox}

\begin{takeawaybox}
Merge sort suits linked lists because split and merge are pointer
operations needing no random access or extra array. It is the default
answer to ``sort a list'' and reuses the middle-finder and merge routines
seen elsewhere.
\end{takeawaybox}
